% !TEX root = task-1-15-3.tex
% Editable task bank. Append new BankTask blocks at the end of the document.
\documentclass[12pt,a4paper]{article}
\usepackage{physics-textbook}
\usepackage{graphicx,needspace}
\newif\ifShowSolutions
\ifdefined\TutorEdition
  \ShowSolutionstrue
\else
  \ShowSolutionsfalse % Change to \ShowSolutionstrue to display solutions.
\fi
\newcounter{banktask}
\newcommand{\BankTask}[3]{%
  \refstepcounter{banktask}%
  \Needspace{35mm}\par\medskip\noindent\begin{minipage}{\linewidth}%
  \noindent\textbf{\thebanktask.}\quad #2\par
  \ifShowSolutions
    #3\par
    \ifdefined\TutorEdition
      \ifcsname TutorSource#1\endcsname
        \smallskip{\small\csname TutorSource#1\endcsname\par}%
      \fi
    \fi
  \fi
  \end{minipage}\par\medskip
}
\begin{document}
\BookSubsection{1.15.3}{Механическая энергия}


% Source sdamgia-519: https://phys-ege.sdamgia.ru/problem?id=519
\BankTask{sdamgia-519}{
Ведро массой $m$ поднимают на веревке вертикально вверх с постоянной скоростью. Когда ведро поднимется на высоту $h$ относительно первоначального положения, его полная механическая энергия в неподвижной системе отсчета
 
$1$) увеличится на величину $mgh$
$2$) уменьшится на величину $mgh$ 
$3$) не изменится
$4$) изменится на неизвестную величину, так как не задана скорость
}{
\smallskip\noindent\textbf{Решение}\par
Поскольку ведро поднимают с постоянной скоростью, его кинетическая энергия не изменяется. Таким образом, полная механическая энергия увеличивается на величину приращение потенциальной энергии, то есть на величину $mgh$.
Ответ: $1$.
}

% Source sdamgia-520: https://phys-ege.sdamgia.ru/problem?id=520
\BankTask{sdamgia-520}{
Ведро массой $m$ опускают на веревке вертикально вниз с постоянной скоростью. Когда ведро опустится на расстояние $h$ относительно первоначального положения, его полная механическая энергия
 
$1$) увеличится на величину $mgh$
$2$) уменьшится на величину $mgh$
$3$) не изменится
$4$) будет неизвестна, так как не задана скорость
}{
\smallskip\noindent\textbf{Решение}\par
Поскольку ведро опускают с постоянной скоростью, его кинетическая энергия не изменяется. Таким образом, полная механическая энергия уменьшается на величину убыли потенциальной энергии, то есть на величину $mgh$.
Ответ: $2$.
}

% Source sdamgia-522: https://phys-ege.sdamgia.ru/problem?id=522
\BankTask{sdamgia-522}{
Мяч массой $m$ бросают вертикально вверх. Если пренебречь сопротивлением воздуха, то при подъеме мяча на высоту $h$ относительно первоначального положения, его полная механическая энергия
 
$1$) увеличится на величину $mgh$
$2$) уменьшится на величину $mgh$
$3$) не изменится
$4$) будет неизвестна, так как не задана начальная скорость
}{
\smallskip\noindent\textbf{Решение}\par
Поскольку сопротивлением воздуха можно пренебречь, для мяча выполняется закон сохранения полной механической энергии. Таким образом, при подъеме на высоту $h$ относительно первоначального положения, его полная механическая энергия не изменится.
Ответ: $3$.
}

% Source sdamgia-39916: https://phys-ege.sdamgia.ru/problem?id=39916
\BankTask{sdamgia-39916}{
Шарик массой $200\,\text{г}$ падает с высоты $10\,\text{м}$ без начальной скорости. К моменту падения на землю потеря полной механической энергии этого шарика из-за действия силы сопротивления воздуха составила $10\%$. Чему равна кинетическая энергия шарика перед ударом о землю?
}{
\smallskip\noindent\textbf{Решение}\par
В верхней точке тело обладало потенциальной энергией $E_{\text{п}}=mgh.$ Так как начальная скорость шарика была равна $0$, то полная механическая энергия равна его потенциальной энергии. В результате потерь из-за работы силы сопротивления воздуха, кинетическая энергия составила $90\%$ от полной энергии. Тогда получаем:


$E_{\text{к}}=0{,}9mgh=0{,}9\cdot0{,}2\cdot10\cdot10=18\,\text{Дж}.$

Ответ: $18$.
}

% Source sdamgia-29022: https://phys-ege.sdamgia.ru/problem?id=29022
\BankTask{sdamgia-29022}{
Камень брошен с горизонтальной поверхности под углом $30^\circ$ к горизонту. Найдите отношение кинетической энергии камня в верхней точке траектории к его начальной кинетической энергии. Сопротивлением воздуха можно пренебречь.
}{
\smallskip\noindent\textbf{Решение}\par
В момент броска камня его кинетическая энергия равна $E_1=\frac{mv_0^2}{2}.$ В верхней точке траектории вертикальная составляющая скорости равна $0$, а горизонтальная равна $v=v_x=v_{0x}=v_0\cos\alpha,$ так как $g_x=0.$ Тогда кинетическая энергия камня в верхней точке траектории равна: 


$E_2=\frac{mv_0^2\cos^2\alpha}{2}.$

Искомое отношение: 


$\frac{E_2}{E_1}=\cos^2\alpha=\frac34=0{,}75.$

Ответ: $0{,}75$.
}

% Source sdamgia-538: https://phys-ege.sdamgia.ru/problem?id=538
\BankTask{sdamgia-538}{
Тело массой $1\,\text{кг}$, брошенное с уровня земли вертикально вверх, упало обратно. Перед ударом о землю оно имело кинетическую энергию $200\,\text{Дж}$. С какой скоростью тело было брошено вверх? Сопротивлением воздуха пренебречь.
}{
\smallskip\noindent\textbf{Решение}\par
Поскольку сопротивлением воздуха можно пренебречь, для тела выполняется закон сохранения полной механической энергии. Значит, кинетическая энергия тела сразу после броска совпадает с кинетической энергией перед ударом о землю. Отсюда находим, с какой скоростью тело было брошено вверх: 


$v=\sqrt{\frac{2E}{m}}=\sqrt{\frac{2\cdot200\,\text{Дж}}{1\,\text{кг}}}=20\,\text{м/с}.$

Ответ: $20$.
}

% Source sdamgia-36749: https://phys-ege.sdamgia.ru/problem?id=36749
\BankTask{sdamgia-36749}{
Шарик, подвешенный на длинной легкой нерастяжимой нити, совершает колебания. Максимальная потенциальная энергия шарика в поле тяжести, если считать ее равной нулю в положении равновесия, равна $0{,}8\,\text{Дж}$. Максимальная скорость шарика в процессе колебаний равна $2\,\text{м/с}$. Чему равна масса шарика? Сопротивлением воздуха пренебречь.
}{
\smallskip\noindent\textbf{Решение}\par
Для математического маятника максимальная потенциальная энергия шарика будет в крайней точке, максимальная скорость — в точке равновесия. Используя закон сохранения энергии $E_{\text{п max}}=E_{\text{к max}}=\frac{mv_{\max}^2}{2},$ получаем:


$m=\frac{2E_{\text{п max}}}{v_{\max}^2}=\frac{2\cdot0{,}8}{4}=0{,}4\,\text{кг}.$

Ответ: $0{,}4$.
}

% Source sdamgia-2941: https://phys-ege.sdamgia.ru/problem?id=2941
\BankTask{sdamgia-2941}{
Снаряд массой $4\,\text{кг}$, летящий со скоростью $400\,\text{м/с}$, разрывается на две равные части, одна из которых летит в направлении движения снаряда, а другая — в противоположную сторону. В момент разрыва суммарная кинетическая энергия осколков увеличилась на величину $\Delta E.$ Скорость осколка, летящего по направлению движения снаряда, равна $900\,\text{м/с}$. Найдите $\Delta E.$
}{
\smallskip\noindent\textbf{Решение}\par
Введем обозначения: $2m$ — масса снаряда до взрыва; $v_0$ — модуль скорости снаряда до взрыва; $v_1$ — модуль скорости осколка, летящего вперед; $v_2$ — модуль скорости осколка, летящего назад. 
Система уравнений для решения задачи:


$\begin{cases}2mv_0=mv_1-mv_2&\text{закон сохранения импульса},\\ \frac{2mv_0^2}{2}+\Delta E=\frac{mv_1^2}{2}+\frac{mv_2^2}{2}&\text{закон сохранения энергии}.\end{cases}$

Выразим $v_2$ из первого уравнения: $v_2=v_1-2v_0$ и подставим во второе уравнение. Получим: 


$v_1^2-2v_0v_1+v_0^2-\frac{\Delta E}{m}=0.$
Отсюда следует: 


$\Delta E=m(v_1-v_0)^2=2\,\text{кг}\cdot(900\,\text{м/с}-400\,\text{м/с})^2=0{,}5\,\text{МДж}.$

Ответ: $\Delta E=0{,}5\,\text{МДж}.$
}

% Source sdamgia-2945: https://phys-ege.sdamgia.ru/problem?id=2945
\BankTask{sdamgia-2945}{
Кусок пластилина сталкивается со скользящим навстречу по горизонтальной поверхности стола бруском и прилипает к нему. Скорости пластилина и бруска перед ударом направлены взаимно противоположно и равны $v_{\text{пл}}=15\,\text{м/с}$ и $v_{\text{бр}}=5\,\text{м/с}.$ Масса бруска в $4$ раза больше массы пластилина. Коэффициент трения скольжения между бруском и столом $\mu=0{,}17.$ На какое расстояние переместятся слипшиеся брусок с пластилином к моменту, когда их скорость уменьшится в $2$ раза?
}{
\smallskip\noindent\textbf{Решение}\par
Пусть $m$ — масса куска пластилина, $M$ — масса бруска, $v_0$ — начальная скорость бруска с пластилином после взаимодействия. Согласно закону сохранения импульса: $Mv_{\text{бр}}-mv_{\text{пл}}=(M+m)v_0.$
Так как $M=4m$ и $v_{\text{бр}}=\frac13v_{\text{пл}},$ то $4m\frac13v_{\text{пл}}-mv_{\text{пл}}=5mv_0$ и $v_0=\frac1{15}v_{\text{пл}}.$
По условию конечная скорость бруска с пластилином $v=0{,}5v_0.$ 
По закону сохранения и изменения механической энергии: 


$\frac{(M+m)v_0^2}{2}=\frac{(M+m)v^2}{2}+\mu(M+m)gS,$
 откуда: 


$\frac{5m(\frac1{15}v_{\text{пл}})^2}{2}=\frac{5m(0{,}5\frac1{15}v_{\text{пл}})^2}{2}+5m\mu gS,$
 


$\frac{v_{\text{пл}}^2}{2\cdot15^2}-\frac{0{,}25v_{\text{пл}}^2}{2\cdot15^2}=\mu gS,$
 


$S=\frac{3}{8\cdot15^2}\frac{v_{\text{пл}}^2}{\mu g}=\frac{3\cdot15^2}{8\cdot15^2\cdot0{,}17\cdot10}\approx0{,}22\,\text{м}.$

Ответ: $S\approx0{,}22\,\text{м}.$
}

% Source sdamgia-42048: https://phys-ege.sdamgia.ru/problem?id=42048
\BankTask{sdamgia-42048}{
Шарик массой $M=10\,\text{кг}$ скользит по гладкому горизонтальному столу со скоростью $V_1=5\,\text{м/с}$ и испытывает лобовое столкновение с покоящимся шариком массой $m=2\,\text{кг}$. После этого скорость тяжёлого шарика остаётся направленной вдоль той же прямой и уменьшается до значения $V_2=4\,\text{м/с}$. Какое количество теплоты $Q$ выделяется в процессе столкновения шариков?
}{
\smallskip\noindent\textbf{Решение}\par
$1$. Поскольку стол гладкий, в процессе столкновения шариков сохраняется проекция импульса системы на прямую, вдоль которой движутся шарики: $MV_1=MV_2+mv_2.$ Отсюда находим скорость второго шарика после столкновения: $v_2=\frac Mm(V_1-V_2)=5\,\text{м/с}.$
$2$. Из закона изменения механической энергии системы получаем, что разность кинетических энергий системы перед столкновением и после него равна искомому количеству теплоты: $E_{k1}-E_{k2}=Q.$ Здесь 


$E_{k1}=\frac{MV_1^2}{2}=10\cdot\frac{25}{2}=125\,\text{Дж},$
$E_{k2}=\frac{MV_2^2}{2}+\frac{mv_2^2}{2}=10\cdot\frac{16}{2}+2\cdot\frac{25}{2}=105\,\text{Дж}.$
$3$. Таким образом, $Q=125-105=20\,\text{Дж}.$
 
Ответ: $Q$ = $20\,\text{Дж}$.
}

\end{document}
